%------------------------------------------------------------------------------%
\documentclass[fleqn,utf8,aspectratio=169,12pt]{beamer}

\usepackage{integracao_e_series}

\title{Usos da Séries de Taylor -- Identidade de Euler}

%------------------------------------------------------------------------------%
\begin{document}

\addtitle

%------------------------------------------------------------------------------%
\section{Identidade de Euler}

%------------------------------------------------------------------------------%
\begin{frame}{Identidade de Euler}
  Com os devidos cuidados todos os resultados sobre sequências e séries\\
  valem para os \emph{números complexos}
  \vspace{\baselineskip}
  \[
    z = a + bi \qquad \text{onde} \qquad a, b \in \R \qquad {\color{blue}i=\sqrt{-1}}
  \]
  \pause
  \vspace{0.5\baselineskip}
  \[
    i^2 = -1 \qquad i^3 = -i \qquad i^4 = 1 \qquad i^5 = i \qquad \cdots
  \]
\end{frame}

%------------------------------------------------------------------------------%
\begin{frame}{Identidade de Euler}
Definições Formais para $z\in\mathbb{C}$
\pause
\[
  \sin(z) = z - \frac{z^3}{3!} + \frac{z^5}{5!} + \frac{z^7}{7!} - \cdots
\]
\pause
\[
  \cos(z) = 1 - \frac{z^2}{2!} + \frac{z^4}{4!} + \frac{z^8}{8!} - \cdots
\]
\pause
\[
  e^z = 1 + z + \frac{z^2}{2!} + \frac{z^3}{3!} + \frac{z^4}{4!} +  \cdots
\]
\end{frame}

%------------------------------------------------------------------------------%
\begin{frame}{Identidade de Euler}
Escolhendo $z=i\theta$
\begin{align*}
  \action<+->{e^{i\theta}}
  \action<+->{& = 1 + (i\theta)
                    + \frac{(i\theta)^2}{2!}
                    + \frac{(i\theta)^3}{3!}
                    + \frac{(i\theta)^4}{4!}
                    + \frac{(i\theta)^5}{5!}
                    + \frac{(i\theta)^6}{6!}
                    + \cdots \\[2mm]}
  \action<+->{& = 1 +  i\theta
                    + \frac{i^2\theta^2}{2!}
                    + \frac{i^3\theta^3}{3!}
                    + \frac{i^4\theta^4}{4!}
                    + \frac{i^5\theta^5}{5!}
                    + \frac{i^6\theta^6}{6!}
                    + \cdots \\[2mm]}
  \action<+->{& = 1 +  i\theta
                    - \frac{\theta^2}{2!}
                    - i\frac{\theta^3}{3!}
                    + \frac{\theta^4}{4!}
                    + i\frac{\theta^5}{5!}
                    - \frac{\theta^6}{6!}
                    + \cdots \\[2mm]}
  \action<+->{& = \left( 1 - \frac{\theta^2}{2!}
                           + \frac{\theta^4}{4!}
                           - \frac{\theta^6}{6!}
                           + \cdots
                  \right)
                + i \left( \theta - \frac{\theta^3}{3!}
                                  + \frac{\theta^5}{5!}
                                  - \cdots
                    \right) \\[2mm]}
  \action<+->{& = \cos\theta + i\sin\theta}
\end{align*}
\end{frame}

%------------------------------------------------------------------------------%
\begin{frame}{Identidade de Euler}
\[
  e^{i\theta} = \cos\theta + i\sin\theta
\]
\pause
``Todas as Constantes da Matemática''
\[
  e^{i\pi} \;=\; \cos\pi + i\sin\pi \;=\; -1 + i0 \;=\; -1
\]
\[
  e^{i\pi} + 1 = 0
\]
\end{frame}

%------------------------------------------------------------------------------%
\begin{frame}{\contentsname}
  \tableofcontents
\end{frame}

%------------------------------------------------------------------------------%
\end{document}
%------------------------------------------------------------------------------%
